\chapter{Título del capítulo}
\label{cap:capitulo-01}

\begin{objetivos}
% Escribe aquí los objetivos del capítulo o elimina este bloque.
\end{objetivos}

\section{Primera sección}
% Escribe aquí tus notas. Puedes dividir el contenido con subsecciones.

\subsection{Subsección}
% Añade ejemplos, figuras, tablas y referencias según las necesites.

% Bloques disponibles en notas.cls: definition, theorem, lemma, proposition,
% corollary, example, exercise, problem, remark y notation.
% Las figuras se pueden guardar en assets/ y añadir con \includegraphics.
